\documentclass[11pt,a4paper]{article}
\input{T0_preamble_local_De.tex}

\title{\textbf{Dreieck-Matrix-Reduktionstheorem}\\
\large Eine geometrische Brückenstruktur zwischen FFGFT und verwandten Formulierungen\\[2mm]
\normalsize Dokument 206 der Reihe \emph{Fundamental Fractal Geometric Field Theory}}
\author{Johann Pascher\\
\small Independent Researcher, Oberösterreich, Austria\\[1mm]
\small ORCID: \href{https://orcid.org/0009-0000-6518-4064}{0009-0000-6518-4064}\\
\small \href{mailto:johann.pascher@gmail.com}{johann.pascher@gmail.com}}
\date{8.\,Mai 2026}

\begin{document}
\maketitle

\begin{abstract}
\noindent
Wir formulieren und beweisen das \emph{Dreieck-Matrix-Reduktionstheorem}: jede konsistente geometrische Beschreibung physikalischer Strukturen, die mit FFGFT in dasselbe Skalenregime fällt, lässt sich auf zwei äquivalente Darstellungen abbilden — eine Z\textsubscript{3}-Dreiecksverbindung und eine 3\,$\times$\,3-Matrix-Algebra. Die Brückenrelation zwischen diesen Darstellungen ist \(\det(A) = 1 - \xi\), woraus die fraktale Dimension \(D_f = 3 - \xi\) als spektrale Größe in linearer Ordnung folgt.

Wir testen das Theorem an sechs konkreten Brücken zu Formulierungen anderer Forscher: Austin (\(D = 3 + \varepsilon\)), Tekermen (UIFT-Offenheit), Phillips (Self-Dual C-25), Porter (\(\Lambda\)-linear), Chekanov-Hakan (trigonometrische Spektren) und Moseley (ITR \(\nu_M\)). Fünf Brücken bestätigen das Theorem algebraisch oder strukturell; eine (Moseley) zeigt eine Regime-Trennung, die kein Widerspruch, sondern eine Klärung der Geltungsbereiche ist.

Wir kontrollieren explizit die Versuchung der Numerologie: ein scheinbarer Treffer für \(\Omega_\Lambda = 0{,}6889\) wird durch einen Filter-Test als nicht-theoretisch identifiziert und nicht als Beweis verbucht. Wir ergänzen einen methodischen Abschnitt zur Zirkularität kosmologischer Vergleichsdaten — \(\Omega_\Lambda\) ist innerhalb des \(\Lambda\)CDM-Rahmens (Dunkle Materie, Standard-Inflation, FRW-Ausdehnung) bestimmt und ist keine ohne Weiteres modellunabhängige Beobachtungsgröße.

Die ehrliche Bilanz: Das Reduktionstheorem ist robust für intrinsische Theorie-Größen. Kosmologische Beobachtungsgrößen wie \(\Omega_\Lambda\) bedürfen einer FFGFT-eigenen Auswertungspipeline auf Rohdaten-Ebene, bevor ein direkter quantitativer Vergleich sinnvoll wird.

Alle Beweise sind reproduzierbar in dreizehn beigefügten Python-Skripten ausgeführt.
\end{abstract}

\tableofcontents

\clearpage
\input{dok206_dreieck_matrix_reduktion_De_ch.tex}

\end{document}
